% Part: incompleteness
% Chapter: arithmetization-syntax
% Section: proofs-in-lk

\documentclass[../../../include/open-logic-section]{subfiles}

\begin{document}

\olfileid[hi]{inc}{art}{plk}
\olsection{$\Log{LK}$ में \usetoken{P}{derivation}}

\begin{explain}
!!{derivation} का अंकगणितीकरण करने के लिए हमें !!{derivation} को संख्याओं के रूप में
निरूपित करना होगा। चूँकि !!{derivation} ऐसे अनुक्रमों के वृक्ष हैं जिनमें
प्रत्येक अनुमान के साथ एक नामपत्र भी होता है, पुनरावर्ती निरूपण सबसे स्वाभाविक
उपाय है: हम किसी !!{derivation} को ऐसे टपल के रूप में निरूपित करते हैं जिसके
घटक अंतिम अनुक्रम, नामपत्र और अंतिम अनुमान के आधारवाक्यों तक पहुँचाने वाली
उप-!!{derivation} के निरूपण हैं।
\end{explain}

\begin{defn}
यदि $\Gamma$, !!{sentence} का परिमित अनुक्रम है, $\Gamma =
\tuple{!A_1, \dots, !A_n}$, तो $\Gn{\Gamma} = \tuple{\Gn{!A_1},
  \dots, \Gn{!A_n}}$।

यदि $\Gamma \Sequent \Delta$ एक अनुक्रम है, तो $\Gamma \Sequent \Delta$ की
ग्योडेल (G\"odel) संख्या है
\[
\Gn{\Gamma \Sequent \Delta} = \tuple{\Gn{\Gamma}, \Gn{\Delta}}
\]

यदि $\pi$, $\Log{LK}$ में !!a{derivation} है, तो $\Gn{\pi}$ निम्न प्रकार
परिभाषित है:
\begin{enumerate}
\item यदि $\pi$ केवल आरंभिक अनुक्रम $\Gamma \Sequent \Delta$ से बनी है, तो
  $\Gn{\pi}$ है
  \[
    \tuple{0, \Gn{\Gamma \Sequent \Delta}}.
  \]
\item यदि $\pi$ एक या दो आधारवाक्यों वाले अनुमान पर समाप्त होती है, उसका
  निष्कर्ष $\Gamma \Sequent \Delta$ है, और $\pi_1$ तथा $\pi_2$ अंतिम अनुमान
  के आधारवाक्यों पर समाप्त होने वाले तात्कालिक उपप्रमाण हैं, तो $\Gn{\pi}$
  क्रमशः है
  \begin{align*}
    & \tuple{1, \Gn{\pi_1}, \Gn{\Gamma \Sequent \Delta}, k} \text{ अथवा}\\ 
    & \tuple{2, \Gn{\pi_1}, \Gn{\pi_2}, \Gn{\Gamma \Sequent \Delta}, k}, 
  \end{align*}
  जहाँ अंतिम अनुमान में प्रयुक्त नियम के अनुसार $k$ निम्न सारणी से मिलता है:

  \begin{tabular}{lccccccc}
    \text{नियम:} & \LeftR{\Weakening} & \RightR{\Weakening} &
    \LeftR{\Contraction} & \RightR{\Contraction} &
    \LeftR{\Exchange} & \RightR{\Exchange} \\
    $k$: & 1 & 2 & 3 & 4 & 5 & 6 \\[2ex]
    \text{नियम:} & \LeftR{\lnot} & \RightR{\lnot} &
    \LeftR{\land} & \RightR{\land} & 
    \LeftR{\lor} & \RightR{\lor} \\
    $k$: & 7 & 8 & 9 & 10 & 11 & 12 \\[2ex]
    \text{नियम:} & \LeftR{\lif} & \RightR{\lif} &
    \LeftR{\lforall} & \RightR{\lforall} &
    \LeftR{\lexists} & \RightR{\lexists} \\
    $k$: & 13 & 14 & 15 & 16 & 17 & 18 \\[2ex]
    \text{नियम:} & \Cut & = \\
    $k$: & 19 & 20
  \end{tabular}
\end{enumerate}
\end{defn}

\begin{ex}
  इस बहुत सरल !!{derivation} पर विचार करें:
  \begin{prooftree}
    \Axiom$!A \fCenter !A$
    \RightLabel{\LeftR{\land}}
    \UnaryInf$!A \land !B \fCenter !A$
    \RightLabel{\RightR{\lif}}
    \UnaryInf$\fCenter (!A \land !B) \lif !A$
  \end{prooftree}
  आरंभिक अनुक्रम की ग्योडेल (G\"odel) संख्या $p_0 = \tuple{0,
    \Gn{!A \Sequent !A}}$ होगी।~$\LeftR{\land}$ के निष्कर्ष पर समाप्त होने
  वाली !!{derivation} की ग्योडेल (G\"odel) संख्या $p_1 = \tuple{1, p_0,
    \Gn{!A \land !B \Sequent !A}, 9}$ होगी ($1$ इसलिए कि $\LeftR{\land}$ का
    एक आधारवाक्य है, $!A \land !B \Sequent !A$ निष्कर्ष की ग्योडेल (G\"odel)
     संख्या है, और $9$,~$\LeftR{\land}$ को कूटित करने वाली संख्या
    है)। तब संपूर्ण !!{derivation} की ग्योडेल (G\"odel) संख्या
    $\tuple{1, p_1, \Gn{\Sequent (!A \land !B) \lif !A)}, 14}$ है, अर्थात्
  \[
  \tuple{1, \tuple{1, \tuple{0, \Gn{!A \Sequent !A)}}, \Gn{!A \land !B \Sequent !A}, 9},
    \Gn{\Sequent (!A \land !B) \lif !A}, 14}.
  \]
\end{ex}

\begin{explain}
!!{derivation} का निरूपण तय कर लेने पर हमें यह भी दिखाना होगा कि ऐसी
!!{derivation} पर आदिम पुनरावर्ती ढंग से क्रियाएँ की जा सकती हैं और उनके
आवश्यक गुण तथा संबंध भी उसी ढंग से व्यक्त किए जा सकते हैं। कुछ संक्रियाएँ
सरल हैं: उदाहरणतः किसी !!a{derivation} की ग्योडेल (G\"odel) संख्या~$p$ दिए
जाने पर $\fn{EndSeq}(p) = (p)_{(p)_0+1}$ उसके अंतिम अनुक्रम की ग्योडेल (G\"odel)
 संख्या और $\fn{LastRule}(p) = (p)_{(p)_0+2}$ उसके अंतिम नियम का कूट
देता है। गुण $\fn{Sequent}(s)$, जिसे
\[
  \len{s} = 2 \land \bforall{i<\len{(s)_0} + \len{(s)_1}}{\fn{Sent}(((s)_0 \concat (s)_1)_i)}
\]
से परिभाषित किया गया है,~$s$ के लिए तभी लागू होता है जब $s$, !!{sentence} से
बने किसी अनुक्रम की ग्योडेल (G\"odel) संख्या हो। कुछ संक्रियाएँ कहीं अधिक
कठिन हैं। हम कम-से-कम यह रूपरेखा देंगे कि उन्हें कैसे किया जाए। लक्ष्य यह
दिखाना है कि ``$\pi$,~$\Gamma$ से~$!A$ की !!a{derivation} है'' संबंध, $\pi$
और~$!A$ की ग्योडेल (G\"odel) संख्याओं का आदिम पुनरावर्ती संबंध है।
\end{explain}

\begin{prop}\ollabel{prop:followsby}
  गुण $\fn{Correct}(p)$, जो तभी लागू होता है जब ग्योडेल (G\"odel) संख्या~$p$
  वाली !!{derivation}~$\pi$ का अंतिम अनुमान सही हो, आदिम पुनरावर्ती है।
\end{prop}

\begin{proof}
  $\Gamma \Sequent \Delta$ आरंभिक अनुक्रम है यदि या तो कोई !!a{sentence}~$!A$
  ऐसा है कि $\Gamma \Sequent \Delta$, $!A \Sequent !A$ है, अथवा कोई पद~$t$
  ऐसा है कि $\Gamma \Sequent \Delta$, $\emptyset \Sequent \eq[t][t]$ है।
  ग्योडेल (G\"odel) संख्याओं के पदों में $\fn{InitSeq}(s)$ लागू होता है यदि
  और केवल यदि
  \begin{align*}
  \bexists{x < s}{} (\fn{Sent}(x) & \land
  s = \tuple{\tuple{x},\tuple{x}})
  \lor {}\\
  \bexists{t<s}{} (\fn{Term}(t) & \land
  s = \tuple{0, \tuple{\Gn{{\eq}(} \concat t \concat \Gn{,} \concat t \concat \Gn{)}}}).
  \end{align*}

  हमें यह भी दिखाना है कि प्रत्येक अनुमान-नियम~$R$ के लिए संबंध
  $\fn{FollowsBy}_R(p)$ आदिम पुनरावर्ती है, जहाँ $\fn{FollowsBy}_R(p)$ तभी
  लागू होता है जब $p$, !!{derivation}~$\pi$ की ग्योडेल (G\"odel) संख्या हो और
  ~$\pi$ का अंतिम अनुक्रम,~$\pi$ की तात्कालिक उप-!!{derivation} से~$R$ के सही
  अनुप्रयोग द्वारा फलित हो।

  एक सरल स्थिति \RightR{\land} नियम की है। यदि $\pi$ किसी सही
  $\RightR{\land}$ अनुमान पर समाप्त होती है, तो वह ऐसी दिखती है:
  \begin{prooftree}
    \AxiomC{}
    \RightLabel{$\pi_1$}
    \Deduce$\Gamma \fCenter \Delta, !A$
    
    \AxiomC{}
    \RightLabel{$\pi_2$}
    \Deduce$\Gamma \fCenter \Delta, !B$

    \RightLabel{\RightR\land}
    \BinaryInf$\Gamma \fCenter \Delta, !A \land !B$
  \end{prooftree}
  अतः !!{derivation} $\pi$ का अंतिम अनुमान $\RightR{\land}$ का सही अनुप्रयोग
  है यदि और केवल यदि !!{sentence} के ऐसे अनुक्रम $\Gamma$ और $\Delta$ तथा
  ऐसे दो !!{sentence} $!A$ और~$!B$ हैं कि $\pi_1$ का अंतिम अनुक्रम $\Gamma
  \Sequent \Delta, !A$, $\pi_2$ का अंतिम अनुक्रम $\Gamma \Sequent \Delta,
  !B$, और $\pi$ का अंतिम अनुक्रम $\Gamma \Sequent \Delta, !A \land !B$ है।
  हमें इसे केवल ग्योडेल (G\"odel) संख्याओं में अनूदित करना है। यदि $s =
  \Gn{\Gamma \Sequent \Delta}$, तो $(s)_0 = \Gn{\Gamma}$ और $(s)_1 =
  \Gn{\Delta}$। अतः $\fn{FollowsBy}_{\RightR{\land}}(p)$ लागू होता है यदि
  और केवल यदि
\begin{align*}
& \bexists{g < p}{\bexists{d < p}{\bexists{a < p}{\bexists{b < p}{\quad}}}} \\
& \qquad \fn{EndSequent}(p) = 
  \tuple{g, d \concat 
    \tuple{\Gn{(} \concat a \concat \Gn{\land} \concat b \concat \Gn{)}}} \land {} \\
& \qquad \fn{EndSequent}((p)_1) = \tuple{g, d \concat \tuple{a}} \land {}\\
& \qquad \fn{EndSequent}((p)_2) = \tuple{g, d \concat \tuple{b}} \land {}\\
& \qquad (p)_0 = 2 \land \fn{LastRule}(p) = 10.
\end{align*}
अलग-अलग पंक्तियाँ क्रमशः व्यक्त करती हैं कि ``ग्योडेल (G\"odel) संख्या~$g$
वाला अनुक्रम~($\Gamma$), ग्योडेल (G\"odel) संख्या~$d$ वाला अनुक्रम~($\Delta$),
ग्योडेल (G\"odel) संख्या~$a$ वाला !!a{formula}~($!A$), और ग्योडेल (G\"odel)
संख्या~$b$ वाला !!a{formula}~($!B$) है'', ऐसे कि ``$\pi$ का अंतिम अनुक्रम
$\Gamma \Sequent \Delta, !A \land !B$ है'', ``$\pi_1$ का अंतिम अनुक्रम
$\Gamma \Sequent \Delta, !A$ है'', ``$\pi_2$ का अंतिम अनुक्रम $\Gamma
\Sequent \Delta, !B$ है'', और ``$\pi$ की दो तात्कालिक उपव्युत्पत्तियाँ हैं
तथा अंतिम अनुमान-नियम $\RightR\land$ (संख्या~$10$ वाला) है''।

~$\pi$ का अंतिम अनुमान $\RightR\lexists$ का सही अनुप्रयोग है यदि और केवल यदि
ऐसे अनुक्रम $\Gamma$ और $\Delta$, !!a{formula}~$!A$, चर~$x$ और पद~$t$ हैं कि
$\pi$ का अंतिम अनुक्रम $\Gamma \Sequent \Delta, \lexists[x][!A]$ और
$\pi_1$ का अंतिम अनुक्रम $\Gamma \Sequent \Delta, \Subst{!A}{t}{x}$ है। अतः
ग्योडेल (G\"odel) संख्याओं के पदों में $\fn{FollowsBy}_{\RightR\lexists}(p)$
लागू होता है यदि और केवल यदि
\begin{align*}
  & \bexists{g<p}{\bexists{d<p}{\bexists{a<p}{\bexists{x<p}{\bexists{t<p}{\quad}}}}}\\
  & \qquad \fn{EndSequent}(p) = 
  \tuple{
    g,
    d \concat \tuple{\Gn{\lexists} \concat x \concat a}
  } \land {}\\
  & \qquad \fn{EndSequent}((p)_1) = 
  \tuple{
    g,
    d \concat \tuple{\fn{Subst}(a, t, x)}
  } \land {}\\
  & \qquad (p)_0 = 1 \land \fn{LastRule}(p) = 18.
\end{align*}
तब हम $\fn{Correct}(p)$ को इस प्रकार परिभाषित करते हैं:
\begin{multline*}
  \fn{Sequent}(\fn{EndSequent}(p)) \land {}\\ 
  [(\fn{LastRule}(p) = 1 \land 
  \fn{FollowsBy}_{\LeftR\Weakening}(p)) \lor \dots \lor {}\\
  (\fn{LastRule}(p) = 20 \land \fn{FollowsBy}_{\eq}(p)) \lor {}\\
  (p)_0 = 0 \land \fn{InitialSeq}(\fn{EndSequent}(p))]
\end{multline*}
पहली पंक्ति सुनिश्चित करती है कि~$d$ का अंतिम अनुक्रम वास्तव में !!{sentence}
से बना अनुक्रम है। अंतिम पंक्ति उस स्थिति को समेटती है जिसमें $p$ केवल एक
आरंभिक अनुक्रम है।
\end{proof}

\begin{prob}
  \olref[inc][art][plk]{prop:followsby} की भाँति निम्न गुण परिभाषित करें:
  \begin{enumerate}
  \item $\fn{FollowsBy}_{\Cut}(p)$,
  \item $\fn{FollowsBy}_{\LeftR\lif}(p)$,
  \item $\fn{FollowsBy}_{\eq}(p)$,
  \item $\fn{FollowsBy}_{\RightR{\lforall}}(p)$।
  \end{enumerate}
  अंतिम गुण के लिए आपको यह भी दिखाना होगा कि आदिम पुनरावर्ती ढंग से जाँचा
  जा सकता है कि ग्योडेल (G\"odel) संख्या~$p$ वाली !!{derivation} का अंतिम
  अनुमान स्वचर-शर्त पूरी करता है या नहीं; अर्थात् $\RightR{\lforall}$ का
  स्वचर~$a$ अंतिम अनुक्रम में नहीं आता।
  \end{prob}
  

\begin{prop}
  \ollabel{prop:deriv}
  संबंध $\fn{Deriv}(p)$, जो तब लागू होता है जब $p$ किसी सही
  !!{derivation}~$\pi$ की ग्योडेल (G\"odel) संख्या हो, आदिम पुनरावर्ती है।
\end{prop}

\begin{proof}
  !!^a{derivation}~$\pi$ तब सही है जब उसका प्रत्येक अनुमान किसी नियम का सही
  अनुप्रयोग हो; अर्थात् उसकी प्रत्येक उप-!!{derivation} किसी सही अनुमान पर
  समाप्त हो। अतः $\fn{Deriv}(d)$ यदि और केवल यदि
  \[
  \bforall{i<\len{\fn{SubtreeSeq}(p)}}{\fn{Correct}((\fn{SubtreeSeq}(p))_i}.
  \]
\end{proof}


\begin{prop}
मान लें कि $\Gamma$, !!{sentence} का आदिम पुनरावर्ती समुच्चय है। तब यह
व्यक्त करने वाला संबंध $\Prf[\Gamma](x, y)$ कि ``किसी परिमित $\Gamma_0
\subseteq \Gamma$ के लिए $x$,~$\Gamma_0 \Sequent !A$ की
!!a{derivation}~$\pi$ का कूट है और $y$,~$!A$ की ग्योडेल (G\"odel) संख्या
है'', आदिम पुनरावर्ती है।
\end{prop}

\begin{proof}
मान लें कि ``$y \in \Gamma$'' आदिम पुनरावर्ती विधेय~$R_\Gamma(y)$ से दिया
गया है। हमें दिखाना है कि $\Prf[\Gamma](x, y)$ आदिम पुनरावर्ती है, जो तब
लागू होता है जब $y$ किसी वाक्य~$!A$ की ग्योडेल (G\"odel) संख्या हो और $x$,
अंतिम अनुक्रम $\Gamma_0 \Sequent !A$ वाली किसी $\Log{LK}$-!!{derivation} का
कूट हो।

पिछली प्रतिज्ञप्ति के अनुसार गुण $\fn{Deriv}(x)$, जो तभी लागू होता है जब $x$,
$\Log{LK}$ में किसी सही !!{derivation}~$\pi$ का कूट हो, आदिम पुनरावर्ती है।
यदि $x$ ऐसा कूट है, तो $\fn{EndSequent}(x)$,~$\pi$ के अंतिम अनुक्रम का कूट
है; इसलिए $(\fn{EndSequent}(x))_0$ अंतिम अनुक्रम के बाएँ पक्ष का कूट और
$(\fn{EndSequent}(x))_1$ दाएँ पक्ष का कूट है। अतः हम ``~$\pi$ के अंतिम
अनुक्रम का दायाँ पक्ष~$!A$ है'' को $\len{(\fn{EndSequent}(x))_1} = 1 \land
((\fn{EndSequent}(x))_1)_0 = x$ से व्यक्त कर सकते हैं। $\pi$ के अंतिम अनुक्रम
का बायाँ पक्ष निस्संदेह स्वतः परिमित है; हमें केवल यह व्यक्त करना है कि उसमें
प्रत्येक वाक्य~$\Gamma$ में है। इस प्रकार हम $\Prf[\Gamma](x, y)$ को परिभाषित
कर सकते हैं:
\begin{align*}
\Prf[\Gamma](x, y) \defiff {}&
\fn{Deriv}(x) \land {} \\
& \bforall{i <
  \len{(\fn{EndSequent}(x))_0}}{R_\Gamma(((\fn{EndSequent}(x))_0)_i)} \land {}\\
& \len{(\fn{EndSequent}(x))_1} = 1 \land ((\fn{EndSequent}(x))_1)_0 = y.
\end{align*}
\end{proof}

\end{document}
